\section{Scope of the retained comparison statements}
\label{sec:v109-scope}
For clarity, we record the precise scope of the earlier allocation
statements. None of their algebraic maximality assertions is used
as a substitute for the physical estimate in
Theorem~\ref{thm:v109-residual}.

\begin{proposition}[Fixed-window scalar maximality and old-family dominance]
\label{prop:v109-maximal}
Fix a Borel comparison subspace, $0<r<\Delta$, and $K>0$.
Consider nonnegative measurable weights $w_+,w_-$ with
$w_++w_-\leq1$, supported on the respective one-sided window
conditions in \eqref{eq:v109-pool}. Among scalar accepted heights
$Z(t)$ satisfying
\[
 0\leq Z(t)\leq U(t),\qquad
 Z(t)\leq K\int(w_++w_-)p_\vee\,d\theta,
\]
the greatest is $V_{r,K}(t)=\min\{U(t),KP_r(t)\}$.
The allocation is global at the roof and does not couple individual
source sheets to comparison sheets.

For a compatible old fixed family with window width $r_{\rm old}$,
physical overlap $N$ and block cost $\kappa$, the explicit conditions
\[
 0<r\leq r_{\rm old},\qquad \ell_{\rm comparison}<\Delta,
 \qquad K=N\kappa
\]
imply $V^{\rm old}(t)\leq V_{r,N\kappa}(t)$ on the same original
source and comparison space.
\end{proposition}
\begin{proof}
The support and one-copy conditions give $w_++w_-\leq Q_r$.
Hence every $Z$ is bounded by the stated minimum. On $P_r>0$ choose
$w_\pm=V_{r,K}q_\pm/(KP_r)$; use zero weights otherwise.
This is a measurable attaining certificate. The separate source
multiplier is $\beta_{r,K}$ defined in
Section~\ref{sec:v109-shortage}.
For the old family, retain only active source blocks at $t$.
Their comparison supports lie in the $r$-pool by the displayed
width condition. Summing the blockwise budget and using physical
overlap at most $N$ gives $V^{\rm old}\leq N\kappa P_r$;
disjoint original source blocks also give $V^{\rm old}\leq U$.
The minimum proves the assertion, with countable families obtained
by monotone convergence on their common roof versions.
\end{proof}

\begin{proposition}[Strictness relative to roof-independent whole-interval blocks]
\label{prop:v109-strictness}
In the interval-source model, the roofwise certificate class is
strictly larger than the old class of roof-independent angular
source blocks that requires a comparison interval to contain the
translate of the \emph{entire} source interval by $[0,r]$ or
$[-r,0]$. This statement compares certificate classes and does not
assert a Lorentz realization.
\end{proposition}
\begin{proof}
Let $s=1/1000$ and $r=1/10$. On two disjoint angular sets of measure
one give both endpoint traces value one. Let their physical intervals
be $(0,s)$ and $(3s/8,r+7s/8)$, and select only the first as source.
Take $\Delta>r+s/2$ and $K=1$. For $3s/8<t<7s/8$ the second
interval contains $[t,t+r]$, so $P_r(t)\geq U(t)=1$.
Neither physical interval has length $r+s$, which is necessary for
the old translated-whole-source inclusion. Angular subdivision does
not alter these intervals. Thus the old class cannot allocate a
nonzero whole-source block while the roofwise class accepts the
source on the indicated roof interval.
\end{proof}

\subsection{Limit order and the remaining estimate}
The constructions in this article are evaluated first at finite
$m,R,n,k$ with $B,r_B,h_B,K_B$ already fixed. Positive physical sums
are completed before taking the essential roof supremum, the radius
and exact-label suprema, and the collision limsup. Only then is $B$
sent to infinity. The assertion that every positive-length segment
contains a sufficiently small window at a fixed count supplies no
uniform modulus for this order.

The remaining estimate is precisely
\[
 \lim_{B\to\infty}\limsup_{m\to\infty}\sup_{R,n,k}m^2
 \esssup_{t\in\mathcal C_{m,R,n,k}(M)}
 \bigl(\mathsf U_{B,m}^{\rm micro}(t)
                 -K_BP_{r_B,m}^{\rm ng}(t)\bigr)_+=0.
\]
No inference to this estimate is made from bounded wordwise optical
weights, local integrability, a finite diagnostic, or a positive
finite-count incidence. Theorem~\ref{thm:v109-residual} proves the
entire visited-grazing payment in this order; all remaining positive
classes in the displayed expression are still part of the original
pointwise problem.
